\chapter{Inequality and Legitimacy}
\label{ch:inequality-legitimacy}

\epigraph{The comfort of the rich depends upon an abundant supply of the poor.}{Voltaire}

\noindent
Every unequal society faces a legitimation problem: why do the many accept the privilege of the few?
The question is ancient (Hume noted that the rulers are always outnumbered by the ruled, so government must rest ultimately on opinion), but doxonomics gives it new precision.
The previous chapters analyzed specific regimes (selfishness, meritocracy, consumerism) that contribute to legitimation.
This chapter integrates them into a unified framework: the doxonomic theory of inequality legitimation.

The central argument: inequality persists not (primarily) because of coercion but because of narrative.
The narrative infrastructure that legitimates inequality (the network of institutions, media outlets, educational practices, and cultural forms that make the existing distribution of rewards seem natural, inevitable, or deserved) is itself a product of material investment.
The wealthy fund the narratives that legitimate their wealth: a feedback loop that doxonomics is uniquely equipped to analyze.

\section{The Legitimation Problem}
\label{sec:ch24-problem}

\subsection{Why legitimation is necessary}

In any society where the top 1\% holds 30--40\% of the wealth (as in the contemporary US and UK), the distribution could in principle be challenged by majoritarian political action.
That it is not (that democratic majorities consistently fail to enact redistributive policies commensurate with the level of inequality) requires explanation.

Three explanations are commonly offered:
\begin{enumerate}
  \item \textbf{Coercion}: the wealthy use force (police, military, legal system) to suppress redistributive demands.
  This is occasionally true but insufficient as a general explanation in democracies with contested elections and free speech.

  \item \textbf{Institutional capture}: the political system is structured to block redistribution regardless of majority preferences (gerrymandering, lobbying, campaign finance, filibuster).
  This is substantially true and well-documented, but it is only part of the story.

  \item \textbf{Narrative legitimation}: the majority \emph{accepts} the existing distribution as fair, natural, or inevitable, or at least does not find it sufficiently objectionable to prioritize redistribution.
  This is the doxonomic explanation.
\end{enumerate}

All three mechanisms operate simultaneously.
The doxonomic contribution is not to deny coercion or institutional capture but to analyze the third mechanism, the one that makes the first two unnecessary for much of the time.

\subsection{The legitimation function}

\begin{definition}[Legitimation Function]
\label{def:legitimation}
\index{legitimation}
The \emph{legitimation function} $L: \R_{\geq 0} \times \R_{\geq 0} \to [0,1]$ maps the level of inequality $I$ and doxonomic investment $\phi_L$ to the population's acceptance level:
\[
  L(I, \phi_L) = \frac{\phi_L}{\phi_L + \kappa(I - I_0)^+}
\]
where $I_0$ is a baseline inequality level below which legitimation is ``free'' (inequality is accepted without narrative investment) and $\kappa > 0$ is the sensitivity parameter measuring how effectively visible inequality erodes acceptance.
\end{definition}

Properties of the legitimation function:
\begin{itemize}[nosep]
  \item $L = 1$ when $I \leq I_0$: below the baseline, inequality requires no justification.
  \item $L$ decreases as $I$ increases for fixed $\phi_L$: rising inequality erodes legitimacy.
  \item $L$ increases as $\phi_L$ increases for fixed $I$: more narrative investment increases acceptance.
  \item $\lim_{\phi_L \to \infty} L = 1$: with unlimited narrative investment, any inequality can be legitimated.
  \item $\lim_{I \to \infty} L = 0$ for fixed $\phi_L$: at extreme inequality, finite investment cannot maintain legitimacy.
\end{itemize}

\subsection{The escalating cost of legitimation}

\begin{proposition}[Escalating Cost of Legitimation]
\label{prop:escalation}
\index{legitimation cost}
For a fixed acceptance level $L^*$, the required narrative investment is:
\[
  \phi_L^* = \frac{L^*}{1 - L^*} \cdot \kappa(I - I_0)
\]
which grows linearly in inequality.
Highly unequal societies must spend proportionally more on narrative maintenance.
\end{proposition}

\begin{example}
\label{ex:legitimation-cost-us}
Applying the model to the United States, 1970--2020:

\begin{center}
\begin{tabular}{lcccc}
\toprule
& 1970 & 1990 & 2010 & 2020 \\
\midrule
Top 1\% wealth share ($I$) & 22\% & 30\% & 35\% & 38\% \\
Think-tank spending (est., \$B) & 0.1 & 0.5 & 0.8 & 1.0 \\
Meritocratic belief (\%) & 70 & 75 & 72 & 68 \\
\bottomrule
\end{tabular}
\end{center}

Inequality rose by 73\% (from 22\% to 38\% of wealth held by the top 1\%) while think-tank spending rose tenfold.
Meritocratic belief initially rose despite growing inequality (the model predicts this: investment can outpace erosion), but began declining after 2010 as inequality outstripped the narrative infrastructure's capacity to legitimate it.
The 2008 crisis and post-2016 populism may mark the onset of legitimation failure, a point we return to below.
(Illustrative data; actual estimation requires the methods of Chapters~\ref{ch:doxometric-measurement}--\ref{ch:causal-inference}.)
\end{example}

\section{The Legitimation Repertoire}
\label{sec:ch24-repertoire}

The narrative infrastructure of inequality legitimation draws on several interlocking belief systems.
Each provides a distinct justification for the existing distribution.

\subsection{The wealth-virtue association}

The most efficient legitimating narrative is the association of wealth with virtue: the belief that the rich \emph{deserve} their wealth because they are more talented, hardworking, or morally superior.
This is the intersection of the meritocratic regime (Chapter~\ref{ch:meritocracy}) and the selfishness regime (Chapter~\ref{ch:human-nature}): competition is natural, markets are fair, winners earned their position, losers deserve their fate.

\begin{model}[Wealth-Virtue Feedback]
\label{mod:wealth-virtue}
\index{wealth-virtue association}
The wealth-virtue association creates a feedback loop:
\begin{enumerate}[nosep]
  \item Wealthy individuals and corporations fund institutions that promote the wealth-virtue association (think tanks, business schools, media).
  \item These institutions produce narratives that frame wealth as deserved.
  \item The public internalizes the association, reducing political demand for redistribution.
  \item Reduced redistribution preserves the wealth that funds the narrative infrastructure.
\end{enumerate}
The loop is self-sustaining as long as the narrative investment is sufficient to maintain legitimation above the threshold $L^*$.
\end{model}

\subsection{Trickle-down narratives}

The ``trickle-down'' narrative holds that inequality benefits everyone because wealth at the top generates investment, innovation, and employment that eventually reach the bottom.
The narrative has been a staple of supply-side economics since the 1980s and is actively promoted by business organizations and free-market think tanks.

The empirical evidence is weak: \textcite{piketty2014} documents that periods of high inequality are not associated with faster growth; several studies find the opposite relationship.
But the narrative persists because it serves a specific doxonomic function: it converts a moral problem (is inequality just?) into a technical claim (does inequality promote growth?) that can be debated endlessly without ever reaching the normative question.

\subsection{Aspirational narratives}

Rather than justifying the existing distribution, aspirational narratives redirect attention to future possibility: ``You too could be rich.''
Lottery advertising, rags-to-riches media, and ``American Dream'' rhetoric serve this function.
The doxonomic operation: convert a statistical reality (low intergenerational mobility) into a subjective experience (possibility of upward mobility), making the current distribution feel temporary rather than structural.

\begin{proposition}[Prospect of Upward Mobility]
\label{prop:poum}
\index{prospect of upward mobility}
The ``prospect of upward mobility'' (POUM) hypothesis holds that tolerance of inequality reflects the belief that one's own position may improve.
Formally, agent $j$ tolerates inequality if:
\[
  \E_j[U(Y_j^{\text{future}})] > U(Y_j^{\text{current}} + r_j)
\]
where $r_j$ is the redistribution $j$ would receive.
If $j$ overestimates their probability of upward mobility (because aspirational narratives are more prevalent than accurate mobility data), they rationally oppose redistribution even when it would benefit them in expectation.
This is the ``temporarily embarrassed millionaire'' effect.
\end{proposition}

\subsection{Blame narratives}

Blame narratives locate the cause of poverty in the poor themselves: lack of effort, bad decisions, moral failings, cultural deficiency.
These narratives operationalize the meritocratic regime's third proposition ($p_3$: failure reflects personal deficiency) and direct attention away from structural causes (deindustrialization, discrimination, healthcare costs, housing markets).

\begin{definition}[Narrative Shield]
\label{def:narrative-shield}
\index{narrative shield}
A \emph{narrative shield} is a set of beliefs that deflects attention from structural causes of inequality to individual explanations.
Formally, a narrative shield is a mapping:
\[
  \sigma: \{\text{structural explanations for outcome } O\} \to \{\text{individual explanations for outcome } O\}
\]
that replaces ``$O$ is caused by institutional structures $S$'' with ``$O$ is caused by individual characteristics $C$.''
The shield reduces the perceived need for collective action by framing outcomes as matters of personal responsibility.
\end{definition}

Common narrative shields:
\begin{itemize}[nosep]
  \item ``The poor are lazy'' (shields from: labor market structure, wage stagnation, precarious employment).
  \item ``If you wanted health insurance, you should have gotten a better job'' (shields from: employer-based insurance system, pre-existing condition exclusions).
  \item ``College is available to anyone who works hard enough'' (shields from: tuition costs, school quality segregation, information asymmetries).
  \item ``If immigrants can succeed, why can't you?'' (shields from: selection effects in immigration, survivorship bias in success stories).
\end{itemize}

\section{The Political Economy of Legitimation}
\label{sec:ch24-political-economy}

\subsection{Who funds legitimation?}

The legitimation infrastructure is not funded by ``the economy'' in the abstract; it is funded by identifiable actors with identifiable interests.

\begin{example}
\label{ex:legitimation-funders}
The institutional infrastructure of inequality legitimation in the US includes:
\begin{enumerate}[nosep]
  \item \textbf{Corporate-funded think tanks}: Heritage Foundation (\$85M/year), Cato Institute (\$40M), AEI (\$60M), Manhattan Institute (\$25M), and hundreds of smaller organizations in the Atlas Network.
  \item \textbf{Family foundations of the ultra-wealthy}: Koch network (\$400M+/election cycle), Scaife foundations, Bradley Foundation, Olin Foundation (now defunct but highly influential in building the intellectual infrastructure).
  \item \textbf{Industry associations}: US Chamber of Commerce (\$200M+/year in lobbying and communications), National Association of Manufacturers, Business Roundtable.
  \item \textbf{Corporate PR and communications departments}: Fortune 500 companies collectively spend billions on public affairs, issue advertising, and narrative management.
\end{enumerate}
Total estimated annual spending on inequality-legitimating narrative infrastructure: \$1--3 billion in direct think-tank and advocacy spending, plus an unknown (but much larger) amount in corporate communications and media.
\end{example}

\subsection{Demand-side factors}

Legitimation succeeds not only because it is well-funded but because audiences have psychological incentives to accept it.
Just-world belief, system justification theory \autocite{jost2004}, and cognitive dissonance reduction all predict that people will be motivated to accept narratives that justify the existing order, especially if they perceive themselves as benefiting from it (or hope to benefit in the future).

\begin{proposition}[System Justification Premium]
\label{prop:system-justification}
\index{system justification}
System justification theory predicts that, holding narrative investment constant, legitimation is higher among:
\begin{enumerate}[nosep]
  \item those who perceive themselves as benefiting from the current system,
  \item those with higher need for cognitive closure (preference for certainty),
  \item those with lower socioeconomic status (paradoxically, because accepting the system reduces the psychological cost of their position: it is less painful to believe one's poverty is deserved than to believe one is a victim of injustice).
\end{enumerate}
The third prediction (that the disadvantaged sometimes embrace legitimating narratives more than the advantaged) is empirically documented and represents a distinctive contribution of system justification theory to doxonomic analysis \autocite{jost2004}.
\end{proposition}

\section{Legitimation Failure and Crisis}
\label{sec:ch24-failure}

\subsection{When legitimation breaks down}

The model predicts that legitimation fails when inequality grows faster than narrative investment can compensate.
Observable signatures of legitimation failure:

\begin{enumerate}[nosep]
  \item \textbf{Declining trust in institutions}: if the institutions that produce legitimating narratives lose credibility, their output loses effectiveness.
  \item \textbf{Rise of populism}: populist movements (left and right) are characterized by explicit rejection of elite legitimacy: ``the system is rigged.''
  \item \textbf{Increased protest activity}: strikes, demonstrations, and social movements increase when the population no longer accepts the distribution as legitimate.
  \item \textbf{Narrative fragmentation}: the unified legitimating narrative fragments into competing frames, reducing the coherence that sustains common-sense capture.
  \item \textbf{Elite anxiety}: increased spending on private security, gated communities, and ``prepper'' culture among the wealthy signals awareness that legitimation is failing.
\end{enumerate}

\subsection{The 2008--present legitimation crisis}

\begin{example}
\label{ex:legitimation-crisis}
The 2008 financial crisis produced the most visible legitimation failure in Western democracies since the 1930s.

\paragraph{The triggering event.}
Banks that had been celebrated as innovative were revealed to be reckless.
Executives who had been compensated as geniuses were exposed as incompetent or fraudulent.
The meritocratic narrative, that financial success reflects talent, was visibly falsified.

\paragraph{The legitimation response.}
The policy response (bailouts without prosecution, austerity for the public, continued bonuses for executives) further eroded legitimacy by making the double standard visible.
Think tanks and business groups invested heavily in counter-narratives: ``the crisis was caused by government regulation'' (deflecting blame from the financial sector), ``we need growth, not redistribution'' (the trickle-down narrative), ``too much regulation will prevent recovery'' (the deregulation narrative).

\paragraph{Partial legitimation recovery.}
By 2015, meritocratic belief had partially recovered in survey data: the narrative infrastructure proved resilient.
But new cracks had appeared: the Occupy movement (2011), the Sanders and Trump campaigns (2016), and the ``rigged system'' discourse all reflected persistent legitimation failure among significant segments of the population.

\paragraph{The COVID-19 shock.}
The pandemic produced a second legitimation shock: ``essential workers'' earning minimum wage while billionaire wealth increased by trillions made the wealth-virtue association difficult to sustain.
Whether this produces lasting legitimation failure or another partial recovery depends on the relative investment in legitimating vs.\ delegitimating narratives in the coming years.
\end{example}

\subsection{A model of legitimation dynamics}

\begin{model}[Legitimation Dynamics]
\label{mod:legitimation-dynamics}
\index{legitimation dynamics}
\begin{align}
  \frac{dL}{dt} &= \alpha \phi_L (1 - L) - \beta (I - I_0)^+ L - \delta_{\text{crisis}} \cdot \mathbf{1}[\text{crisis at } t] \label{eq:legit-dynamics} \\
  \frac{dI}{dt} &= \gamma_I \cdot r(L) - \mu \label{eq:ineq-dynamics}
\end{align}
where:
\begin{itemize}[nosep]
  \item $\alpha \phi_L (1-L)$: narrative investment increases legitimation (with diminishing returns),
  \item $\beta(I - I_0)^+ L$: excess inequality erodes legitimation,
  \item $\delta_{\text{crisis}}$: a crisis event produces a discrete legitimation shock,
  \item $r(L)$: the rate of return to capital increases with legitimation (high $L$ blocks redistribution and regulation),
  \item $\mu$: a natural compression force (growth, diffusion, demographic change).
\end{itemize}
The system can exhibit:
\begin{enumerate}[nosep]
  \item \emph{Stable legitimation}: $L$ and $I$ reach a self-sustaining equilibrium.
  \item \emph{Vicious spiral}: rising $I$ erodes $L$, but reduced $L$ fails to produce redistribution because institutional capture blocks the political channel, so $I$ continues to rise, further eroding $L$.
  \item \emph{Crisis-triggered regime change}: a legitimation shock ($\delta_{\text{crisis}}$ large enough) pushes $L$ below a threshold, triggering redistributive political action that reduces $I$.
\end{enumerate}
\end{model}

\section{Inequality, Race, and Compound Legitimation}
\label{sec:ch24-race}

Economic inequality intersects with racial and gender hierarchies (Chapter~\ref{ch:race-coloniality-gender}), creating \emph{compound legitimation}: each hierarchy provides narrative resources for legitimating the others.

\begin{proposition}[Compound Legitimation]
\label{prop:compound}
\index{compound legitimation}
When economic inequality correlates with racial hierarchy, each provides narrative cover for the other:
\begin{itemize}[nosep]
  \item Economic inequality is legitimated by racial narratives: ``Those people are poor because of their culture/behavior.''
  \item Racial hierarchy is legitimated by economic narratives: ``If racism were real, how do you explain successful minorities?''
\end{itemize}
The compound effect means that neither inequality nor racial hierarchy can be fully delegitimated without also challenging the other.
Intersectional doxonomic analysis is therefore not a supplement to economic analysis but a requirement for it.
\end{proposition}

Historically, racial division has been a powerful tool for blocking class-based redistribution.
\textcite{du-bois2007} identified the ``psychological wage of whiteness'': the psychic benefit that poor whites derive from racial superiority, which compensates for economic deprivation and reduces demand for redistribution.
This is a doxonomic mechanism: the racial narrative substitutes for material benefit, reducing the political threat that economic inequality would otherwise face.

\section{Comparative Legitimation Regimes}
\label{sec:ch24-comparative}

Different societies legitimate their inequality structures through different narrative repertoires:

\begin{center}
\begin{tabular}{p{2.5cm}p{3cm}p{3.5cm}p{3cm}}
\toprule
Society & Inequality level & Primary legitimation narrative & Counter-narrative strength \\
\midrule
United States & Very high (Gini $\approx 0.39$) & Meritocracy, aspiration & Moderate (populist, progressive) \\[4pt]
United Kingdom & High (Gini $\approx 0.35$) & Tradition, meritocracy & Moderate (labor movement) \\[4pt]
Germany & Moderate (Gini $\approx 0.31$) & Social market economy & Strong (unions, SPD tradition) \\[4pt]
Sweden & Low (Gini $\approx 0.27$) & Social solidarity & Very strong (universal welfare) \\[4pt]
South Africa & Very high (Gini $\approx 0.63$) & Post-apartheid aspiration & Strong but contested \\
\bottomrule
\end{tabular}
\end{center}

The comparison supports the model: higher inequality requires stronger legitimating narratives, and the specific narrative used depends on the available institutional infrastructure.
The US and Sweden are polar cases: the US has high inequality and a strong meritocratic infrastructure; Sweden has low inequality and a strong solidarity infrastructure.
The causal direction runs both ways: institutions produce narratives that sustain inequality levels, and inequality levels demand institutions that produce legitimating narratives.

\section{Notes and References}
\label{sec:ch24-notes}

On inequality and its justification, see \textcite{piketty2014} and \textcite{stiglitz2012}.
The legitimation function draws on \textcite{lukes2005} on the third dimension of power and Gramsci's theory of hegemony.
System justification theory is developed in \textcite{jost2004}.
The wealth-virtue association is analyzed in \textcite{sandel2012}.

On the prospect of upward mobility, see Benabou and Ok (2001).
On the racial dimensions of inequality legitimation, see \textcite{du-bois2007} on the wages of whiteness and \textcite{hochschild2016} on the American Dream's racial contours.
For the political economy of inequality, see \textcite{acemoglu2012} and Hacker and Pierson (2010).

The 2008 legitimation crisis is analyzed in Streeck (2016).
On populism as a symptom of legitimation failure, see M\"{u}ller (2016).
For the Koch network's institutional infrastructure, see \textcite{mayer2016}.

\section{Exercises}

\begin{exercise}
Using the legitimation function (Definition~\ref{def:legitimation}), estimate $\kappa$ for the United States by comparing inequality growth and think-tank spending over 1980--2020.
\begin{enumerate}[nosep]
  \item Gather data on the top 1\% wealth share (as a proxy for $I$) and total think-tank spending (as a proxy for $\phi_L$).
  \item Estimate $I_0$ (the inequality level at which legitimation becomes necessary) from the historical data.
  \item Calibrate $\kappa$ so that the model reproduces observed meritocratic belief levels.
  \item Does the model predict the post-2008 decline in meritocratic belief?
\end{enumerate}
\end{exercise}

\begin{exercise}
Identify three narrative shields (Definition~\ref{def:narrative-shield}) currently active in public discourse about wealth inequality.
For each:
\begin{enumerate}[nosep]
  \item State the structural explanation being shielded.
  \item State the individual explanation being substituted.
  \item Trace the belief production chain: who produces this shield, through what institutions, with what funding?
  \item How would you empirically test whether the shield is effective (i.e., whether exposure to the shield narrative reduces support for structural explanations)?
\end{enumerate}
\end{exercise}

\begin{exercise}
Model a society where legitimation fails ($L < L^*$).
\begin{enumerate}[nosep]
  \item What observable consequences follow from legitimation failure?
  \item Using Model~\ref{mod:legitimation-dynamics}, identify the parameter conditions under which a crisis event triggers lasting legitimation failure vs.\ temporary disruption followed by recovery.
  \item Connect to Chapter~\ref{ch:how-belief-regimes-fracture}: how does legitimation failure relate to belief regime fracture?
\end{enumerate}
\end{exercise}

\begin{exercise}
Apply the compound legitimation framework (Proposition~\ref{prop:compound}) to a specific historical case.
\begin{enumerate}[nosep]
  \item Choose a society where economic and racial inequality are correlated.
  \item Show how economic narratives legitimate racial hierarchy and vice versa.
  \item Identify a historical moment when the compound was challenged (e.g., a cross-racial labor movement) and explain why it succeeded or failed.
\end{enumerate}
\end{exercise}

\begin{exercise}
Construct the comparative legitimation table (Section~\ref{sec:ch24-comparative}) using actual data.
\begin{enumerate}[nosep]
  \item Find Gini coefficients and top-income shares for at least four countries.
  \item Find survey data on meritocratic or just-world belief (World Values Survey, ISSP).
  \item Map the legitimation infrastructure in each country (think tanks, media, policy design).
  \item Does the model's prediction (higher inequality $\to$ more legitimation investment) hold cross-nationally?
\end{enumerate}
\end{exercise}

\begin{exercise}
Analyze the POUM hypothesis (Proposition~\ref{prop:poum}) using mobility data.
\begin{enumerate}[nosep]
  \item Find actual intergenerational mobility rates for the US (e.g., Chetty et al.).
  \item Find survey data on \emph{perceived} mobility (what people think their chances of upward mobility are).
  \item Compute the ``mobility perception gap'': the difference between perceived and actual mobility.
  \item Estimate the doxonomic investment needed to maintain this gap (i.e., to sustain the belief in mobility despite contrary evidence).
\end{enumerate}
\end{exercise}

\begin{exercise}[Computational]
Simulate the legitimation dynamics model (Model~\ref{mod:legitimation-dynamics}):
\begin{enumerate}[nosep]
  \item Set $I_0 = 20$, $I(0) = 25$, $L(0) = 0.7$, $\alpha = 0.1$, $\phi_L = 3$, $\beta = 0.02$, $\gamma_I = 0.5$, $\mu = 0.3$.
  \item Simulate 200 time steps with no crisis ($\delta_{\text{crisis}} = 0$). Plot $L(t)$ and $I(t)$.
  \item At $t = 100$, introduce a crisis ($\delta_{\text{crisis}} = 0.2$). How does the system respond?
  \item Experiment with different $\phi_L$ values. What is the minimum narrative investment needed to prevent legitimation collapse after the crisis?
  \item Compare scenarios where the crisis triggers redistributive policy ($I$ decreases) vs.\ where institutional capture blocks redistribution ($I$ continues rising despite low $L$).
\end{enumerate}
\end{exercise}
